\section{Related Work}
\label{sec:related}

% Mini-outline (one message per paragraph):
% - SL prediction methods: graph/MF link prediction on the SL graph; degree/topology reliance; weak cold-start. (F7: DDGCN/etc. are literature, not our results)
% - Virtual-cell / perturbation models: model the response; closed vocabulary; DL not yet beating linear.
% - Co-dependency & essentiality: DepMap population fitness; essentiality != SL (F4); transcriptome reflects SL.

\paragraph{Computational synthetic-lethality prediction.}
Most SL predictors treat the problem as link prediction over a known SL
interaction graph curated in databases such as SynLethDB
\citep{wang2022synlethdb} and SLKB \citep{gokbag2023slkb}. Matrix-factorization
and graph-neural-network methods (for example the GRSMF, SL2MF, and DDGCN
families surveyed by \citet{wang2022slreview}) propagate observed SL labels
through gene--gene similarity or interaction topology, and the unified benchmark
of \citet{feng2024benchmarking} shows they perform strongly when test pairs share
genes with the training graph. Their common weakness is the same as their
strength: because the signal lives in graph topology and node degree, accuracy
degrades sharply when one or both genes are unseen, and a degree-only probe can
top the easiest split (our Table~\ref{tab:floor}). We use that benchmark and its
difficulty splits as our evaluation harness, but our features come from a gene's
\emph{perturbation response} rather than its position in the SL graph, which is
what gives the approach a path to genes with no curated interactions.

\paragraph{Virtual-cell and perturbation-response models.}
A second line of work models the cell's response to perturbation directly. The
virtual-cell program \citep{bunne2024virtualcell} and the associated challenge
\citep{roohani2025virtualcellchallenge} frame this as learning to predict how a
cell-state distribution shifts under genetic perturbation; STATE
\citep{adduri2025state} is a transformer that predicts such set-level responses,
and single-cell foundation models such as scGPT \citep{cui2024scgpt} and
Tahoe-x1 \citep{gandhi2025tahoex1} pretrain transferable embeddings at scale.
Two limitations matter for our setting. First, these models identify a
perturbation through a fixed gene vocabulary, so out-of-vocabulary genes cannot
be queried without an adapter --- the bottleneck our ESM2-conditioned adapter
\citep{lin2023esm2} targets. Second, careful benchmarking finds that deep
perturbation predictors do not yet consistently beat simple linear baselines
\citep{ahlmanneltze2025linear, vinastorne2025systema}, which is why we frame our
generative framework as a preliminary method and compare it candidly against a
strong dependency-only floor rather than claiming superiority.

\paragraph{Cancer dependency maps and the essentiality--SL distinction.}
A third line characterizes single-gene dependencies from genome-scale CRISPR
screens. The DepMap effort, using CERES/Chronos correction
\citep{meyers2017ceres}, quantifies how much each gene's loss impairs fitness,
and downstream work builds clinically informed and translational dependency maps
and target-prioritization frameworks \citep{pacini2024depmap, shi2024translational},
including deep-learning prediction of dependency from tumor omics
\citep{chiu2021deepdep}. We rely on these dependency scores as features and as a
prediction target for a gene's own essentiality, but we are careful not to
conflate essentiality with synthetic lethality: a broadly essential gene is not
an SL partner, and our decomposition shows that genome-wide co-dependency signal
is dominated by pan-essentiality rather than pair-specific interaction. The
observation closest to our hypothesis is that tumor transcriptomic architecture
reflects SL interactions \citep{haider2025transcriptomic}; we make the predictive
direction explicit and, crucially, generate the transcriptome rather than
requiring it to have been measured.
